% SPDX-License-Identifier: Apache-2.0
% covers: Buffer
\datasheet{Buffer}{The channel as hardware: an elastic buffer of two.}

\dsfacts{\code{lib/parts/src/buffer.rs}}{\code{txhdl\_parts::buffer::Buffer}}%
{\code{//docs:runtime}, \code{//docs:examples}}

\subsection*{Function}

A channel in a simulation is a buffer of two between a sender and a
receiver. \code{Buffer} is that buffer as a netlist, for the places
where two lowered units meet on a channel: a board top, or a design
wired by hand. Its sender side is \code{tx\_data}, \code{tx\_valid} and
\code{tx\_ready}; its receiver side is \code{rx\_data}, \code{rx\_valid}
and \code{rx\_ready}.

\subsection*{Parameters}

\begin{center}\footnotesize
\begin{tabular}{@{}lp{0.7\textwidth}@{}}
\toprule
Parameter & Meaning \\
\midrule
\code{W} & The width of the words the channel holds, in bits. \\
\bottomrule
\end{tabular}
\end{center}

\dstables{Buffer}

\subsection*{Behaviour}

\begin{itemize}
\item \code{tx\_ready} is high when the tail is empty, as the last edge
left it. \code{rx\_valid} is high when the head is full, and
\code{rx\_data} is the head.
\item All three outputs are registers, so there is no combinational
path from one side to the other.
\item At an edge, a take moves the tail into the head, and a push goes
into the head if the head is then empty, else into the tail.
\item With both sides running every cycle, one word passes per cycle.
\end{itemize}

\subsection*{Verification}

\begin{itemize}
\item \code{buffer\_is\_the\_runtime\_channel} in \code{txhdl\_parts}
runs the buffer beside a runtime channel for 200 cycles of pseudorandom
offers and takes, and checks \code{ready}, \code{valid} and the head
after every cycle.
\item \code{ex\_buffer} does the same on a fixed pattern, and the build
simulates \code{Buffer<8>}'s netlist against that run under nvc and
Verilator: \code{//docs:buffer\_sim\_buffer8\_tb\_test} and
\code{//docs:buffer\_vsim\_test}.
\end{itemize}

\subsection*{Used by}

The example \code{ex\_buffer}. A unit of units does not need it, since
its netlist puts a \code{txhdl\_chan} of the same rules between its
children.

\subsection*{Limits}

It holds two words. A consumer that takes in bursts needs \code{Fifo}.
